%======================================================================
\chapter{Correspondence with the formal proof}\label{app:formal}
%======================================================================

The table lists the definitions and results of this volume that carry a ``Formal counterpart''
line, together with the corresponding declarations of the Lean development. All declaration
names start with \lean{HPC.}, except the few Mathlib names marked as such; file names are relative
to the root of the repository. The formal statements are sometimes slightly more general, or are
split differently into lemmas; the table points to the place where the content of each result is
proved. Rows are in the order of the text.

The main theorems of the volume are, in the Lean development (prefix \lean{HPC.} omitted):
\begin{center}
\begin{tabular}{@{}p{0.42\textwidth}p{0.54\textwidth}@{}}
\toprule
Boundary bumping (Theorem~\ref{thm:bump-closed}) & \lean{Jones.exists_continuum_bump_compact},
\lean{Plane.exists_continuum_bump_closed}\\
Moore's non-cut point theorem (Theorems~\ref{thm:noncut-side}, \ref{thm:two-noncut}) &
\lean{Plane.exists_nonCutPoint_of_cut}, \lean{Plane.exists_two_nonCutPoints}\\
Sorgenfrey's theorem (Theorem~\ref{thm:sorgenfrey}) & \lean{Plane.sorgenfrey}\\
Snake-like continua are hereditarily unicoherent and atriodic
(Theorems~\ref{thm:snake-unicoherent}, \ref{thm:snake-atriodic}) &
\lean{Plane.isHereditarilyUnicoherent_of_isSnakeLike}, \lean{Plane.isAtriodic_of_isSnakeLike}\\
Snake-likeness is topological (Theorem~\ref{thm:snake-homeo}) & \lean{BJ59.isSnakeLike_of_homeomorph}\\
Decomposition spaces are metrizable (Theorem~\ref{thm:quotient-metrizable}) &
\lean{Jones.metrizableSpace_quotient}\\
Continuity at a residual set of elements (Theorem~\ref{thm:residual}) & \lean{Jones.residual_lscAt}\\
The open mapping principle (Theorem~\ref{thm:effros-open}) & \lean{Effros.effros_open_mapping}\\
Ungar's corollary (Theorem~\ref{thm:ungar}) & \lean{Effros.effros_homogeneous}\\
\bottomrule
\end{tabular}
\end{center}
All of them are proved in the Lean development without \lean{sorry} and without axioms beyond the
standard ones of Lean and Mathlib.

\begin{longtable}{@{}p{0.40\textwidth}p{0.56\textwidth}@{}}
\toprule
In these notes & In the Lean development\\
\midrule
\endhead
\bottomrule
\endfoot
\multicolumn{2}{@{}l}{\emph{Chapter~\ref{ch:bumping}}}\\
Lemma~\ref{lem:sep-of-not-conn} & \lean{HPC.Plane.exists_separated_of_not_isPreconnected}.\\
Lemma~\ref{lem:side} & \lean{HPC.Plane.subset_side_of_isPreconnected}.\\
Lemma~\ref{lem:union-continua} & \lean{HPC.Plane.isContinuum_union}.\\
A continuum plus one piece (Lemma~\ref{lem:piece}) & \lean{HPC.Plane.isContinuum_union_piece} (Nadler, \emph{Continuum Theory}, 6.3). A version for an arbitrary connected topological space is Theorem~\ref{thm:nadler63}.\\
Lemma~\ref{lem:nested} & \lean{HPC.T14.isPreconnected_sInter_chain}.\\
Lemma~\ref{lem:quasi} & Mathlib's \lean{connectedComponent_eq_iInter_isClopen}, used inside the proofs below.\\
Lemma~\ref{lem:clopen-split} & \lean{HPC.Plane.exists_clopen_split}.\\
Separating two closed sets in a compact set (Theorem~\ref{thm:closed-split}) & \lean{HPC.Plane.exists_closed_split_of_no_connected} (Moore's Theorem~44 in \emph{Foundations of Point Set Theory}). A variant with a continuum in place of the preconnected set, in a compact space, is \lean{HPC.Plane.closed_dichotomy}.\\
Boundary bumping (Theorem~\ref{thm:bump-closed}) & \lean{HPC.Jones.exists_continuum_bump_compact}; for $M$ a continuum in a compact space and $F$ closed, \lean{HPC.Plane.exists_continuum_bump_closed} (Nadler 5.4).\\
Boundary bumping, open form (Corollary~\ref{cor:bump-open}) & \lean{HPC.Plane.exists_continuum_bump_open} (Nadler 5.7).\\
Boundary bumping, metric form (Corollary~\ref{cor:bump-ball}) & \lean{HPC.Plane.exists_subcontinuum_meeting_closedBall}.\\[4pt]
\multicolumn{2}{@{}l}{\emph{Chapter~\ref{ch:irreducible}}}\\
Definition~\ref{def:irreducible} & \lean{HPC.IsIrreducibleBetween} (in \lean{HPC/PseudoArc/Theorem14.lean}).\\
Definition~\ref{def:irr-sets} & \lean{HPC.Plane.IrrFromTo} (Nadler 11.29).\\
Existence of irreducible subcontinua (Theorem~\ref{thm:irr-exists}) & \lean{HPC.Plane.exists_irrFromTo} (Nadler 11.30) and \lean{HPC.exists_irreducible}.\\
Theorem~\ref{thm:indec-irr} & \lean{HPC.Plane.exists_irreducible_of_indecomposable} (the case of Nadler 11.15 used later).\\
Proposition~\ref{prop:irr-homeo} & \lean{HPC.IsIrreducibleBetween.homImage}.\\[4pt]
\multicolumn{2}{@{}l}{\emph{Chapter~\ref{ch:noncut}}}\\
Definition~\ref{def:separation} & \lean{HPC.Plane.IsSeparation}, \lean{HPC.Plane.IsCutPoint} (in \lean{HPC/Plane/NonCut51Gen.lean}).\\
Lemma~\ref{lem:sep-basic} & \lean{HPC.Plane.IsSeparation.disjoint_closure_left}, \lean{HPC.Plane.IsSeparation.isOpen_right}, \lean{HPC.Plane.IsSeparation.subset_or_subset}, \lean{HPC.Plane.exists_separation_of_not_isConnected}.\\
Theorem~\ref{thm:nadler63} & \lean{HPC.Plane.isConnected_union_of_separation} (Nadler 6.3).\\
Lemma~\ref{lem:nadler64} & \lean{HPC.Plane.nadler_6_4}.\\
Moore (Theorem~\ref{thm:noncut-side}) & \lean{HPC.Plane.exists_nonCutPoint_left} and \lean{HPC.Plane.exists_nonCutPoint_of_cut} (Nadler 6.6). Nadler assumes $X$ nondegenerate; this follows from the separation and is not needed.\\
Theorem~\ref{thm:two-noncut} & \lean{HPC.Plane.exists_two_nonCutPoints}.\\
Theorem~\ref{thm:noncut-not-in} & \lean{HPC.Plane.not_subset_of_nonCutPoints} (Nadler 6.7).\\[4pt]
\multicolumn{2}{@{}l}{\emph{Chapter~\ref{ch:triods}}}\\
Definition~\ref{def:triod} & \lean{HPC.Plane.IsTriod}, \lean{HPC.Plane.IsAtriodic} (in \lean{HPC/Plane/Atriodic51Defs.lean}); for a triod in a non-compact space, \lean{HPC.Hagopian.IsTriodSet}.\\
Definition~\ref{def:unicoherent} & \lean{HPC.Plane.IsUnicoherent}, \lean{HPC.Plane.IsHereditarilyUnicoherent}.\\
Definition~\ref{def:esum} & \leansub{HPC.Plane.IsEssentialSum}{2}, \leansub{HPC.Plane.IsEssentialSum}{3}, \lean{HPC.Plane.IsWeakTriod} (in \lean{HPC/Plane/SorgenfreyDefs.lean}); the constructor used in all proofs below is \lean{HPC.Plane.isWeakTriod_of}.\\
Lemma~\ref{lem:esum-basic} & The lemmas in the namespace \leansub{HPC.Plane.IsEssentialSum}{2} (\lean{diff_nonempty}, \lean{closure_diff_subset_left}, \lean{eq_closure_union}, \lean{closure_diff_ne}, \lean{closure_diff_disjoint}, \lean{closure_inter_nonempty}).\\
A triod with a given core (Lemma~\ref{lem:core}) & \lean{HPC.Plane.isTriod_of_core}.\\
Proposition~\ref{prop:triod-weak} & \lean{HPC.Plane.isWeakTriod_of_isTriod} (Nadler 11.25).\\
Corollary~\ref{cor:indec-not-triod} & \lean{HPC.Hagopian.not_isTriod_univ_of_indecomposable} (for the whole space, which is all that is used).\\
Lemma~\ref{lem:weak-not-irr} & \lean{HPC.Plane.not_irreducible_of_isWeakTriod}.\\
Proposition~\ref{prop:uni-triod} & \lean{HPC.Plane.isTriod_iff_isWeakTriod} (Nadler 11.26).\\
Theorem~\ref{thm:diff-connected} & \lean{HPC.Plane.isConnected_diff_of_not_weakTriod} (Nadler 11.27), with the two halves in \lean{HPC.Plane.isConnected_diff_left_of_not_weakTriod}; the continuum $\cl U$ is \lean{HPC.Plane.isContinuum_closure_diff}.\\
Lemma~\ref{lem:1128} & \lean{HPC.Plane.isConnected_diff_of_superset'} (Nadler 11.28, corrected).\\
Proposition~\ref{prop:1131} & \lean{HPC.Plane.nadler_11_31} (Nadler 11.31, without the unused hypothesis $B\subseteq\cl U$).\\
Proposition~\ref{prop:1132} & \lean{HPC.Plane.nadler_11_32} (Nadler 11.32, without the hypothesis $A\ne\cl U$, which follows from the others).\\
Lemma~\ref{lem:1133a} & \lean{HPC.Plane.diff_subset_of_irr}.\\
Proposition~\ref{prop:1133} & \lean{HPC.Plane.nadler_11_33} (Nadler 11.33).\\
Lemma~\ref{lem:sorg-proper} & \lean{HPC.Plane.exists_proper_of_not_irr}.\\
Lemma~\ref{lem:sorg-step} & \lean{HPC.Plane.SorgGood}, \lean{HPC.Plane.sorgGood_base}, \lean{HPC.Plane.sorgGood_step}.\\
Lemma~\ref{lem:sorg-star} & \lean{HPC.Plane.exists_irrFromTo_point} (Nadler 11.34(a)); the body of the proof by contradiction is \lean{HPC.Plane.sorgenfrey_contra}.\\
Sorgenfrey (Theorem~\ref{thm:sorgenfrey}) & \lean{HPC.Plane.sorgenfrey} (Nadler 11.34).\\
Corollary~\ref{cor:irr-iff} & \lean{HPC.Plane.irreducible_iff_not_isTriod} (Nadler 11.35).\\
Corollary~\ref{cor:atriodic-irr} & \lean{HPC.Plane.exists_irreducibleBetween_of_atriodic}.\\[4pt]
\multicolumn{2}{@{}l}{\emph{Chapter~\ref{ch:snake}}}\\
Theorem~\ref{thm:snake-sub} & \lean{HPC.Plane.isSnakeLike_of_subcontinuum}; the interval of links meeting $K$ is \lean{HPC.Plane.meets_between}.\\
Intermediate values along a chain (Lemma~\ref{lem:chain-ivt}) & \lean{HPC.Plane.chain_ivt}.\\
Corollary~\ref{cor:two-sets} & \lean{HPC.Plane.chain_two_sets}.\\
Theorem~\ref{thm:snake-unicoherent} & \lean{HPC.Plane.isUnicoherent_of_isSnakeLike} and \lean{HPC.Plane.isHereditarilyUnicoherent_of_isSnakeLike} (Nadler 12.2).\\
Theorem~\ref{thm:snake-atriodic} & \lean{HPC.Plane.isAtriodic_of_isSnakeLike}, through \lean{HPC.Plane.not_isTriod_of_isSnakeLike} (Nadler 12.4).\\
Proposition~\ref{prop:chain-epsmap} & \lean{HPC.BJ59.exists_epsMap_of_chain} and \lean{HPC.BJ59.exists_epsMap_of_isSnakeLike}.\\
Lemma~\ref{lem:uniform-fibres} & \lean{HPC.Plane.exists_delta_of_fibres_on}.\\
The $\varepsilon$-map criterion (Theorem~\ref{thm:epsmap-criterion}) & \lean{HPC.Plane.isSnakeLike_of_epsMap}; for the whole space, \lean{HPC.Plane.isSnakeLike_univ_of_epsMap}.\\
Snake-likeness is topological (Theorem~\ref{thm:snake-homeo}) & \lean{HPC.BJ59.isSnakeLike_of_homeomorph}.\\
Definition~\ref{def:circchain} & \lean{HPC.Plane.CircChain}, \lean{HPC.Plane.IsCircEpsChain}, \lean{HPC.Plane.IsCircleLike} (in \lean{HPC/Plane/Ingram68.lean}).\\
Proposition~\ref{prop:circle-circlelike} & \lean{HPC.BJ59.isCircleLike_circle}, with the arcs \lean{HPC.BJ59.circArc}.\\
Proposition~\ref{prop:circle-homeo} & \lean{HPC.BJ59.isCircleLike_univ_of_homeomorph}.\\
Theorem~\ref{thm:circle-sub} & \lean{HPC.Plane.isSnakeLike_of_subset_circleLike} (Nadler 12.51).\\[4pt]
\multicolumn{2}{@{}l}{\emph{Chapter~\ref{ch:decomp}}}\\
Definition~\ref{def:usc} & \lean{HPC.Jones.IsUSCDecomp}, \lean{HPC.Jones.LSCAt}, \lean{HPC.Jones.IsContinuousDecomp}, with the element $\elt x$ written \lean{HPC.Jones.cls} (in \lean{HPC/Jones/Decomp.lean}).\\
Lemma~\ref{lem:usc-open} & \lean{HPC.Jones.isOpen_setOf_cls_subset}.\\
Proposition~\ref{prop:closed-usc} & \lean{HPC.Jones.isUSCDecomp_of_isClosed}.\\
Theorem~\ref{thm:quotient-metrizable} & \lean{HPC.Jones.isClosedMap_mk}, \lean{HPC.Jones.t2Space_quotient}, \lean{HPC.Jones.secondCountable_quotient}, \lean{HPC.Jones.metrizableSpace_quotient}. Urysohn's theorem is Mathlib's \lean{TopologicalSpace.metrizableSpace_of_t3_secondCountable}.\\
Theorem~\ref{thm:monotone-preimage} & \lean{HPC.Jones.isContinuum_preimage}.\\
Corollary~\ref{cor:monotone-decomposable} & For monotone maps onto a metric space, \lean{HPC.Hagopian.isDecomposableSpace_of_monotone_image}.\\
Proposition~\ref{prop:homeo-descends} & \lean{HPC.Jones.exists_homeomorph_quotient}.\\
Lemma~\ref{lem:frontier-nd} & \lean{HPC.Jones.dense_compl_frontier_of_isClosed}.\\
Kuratowski, Fort (Theorem~\ref{thm:residual}) & \lean{HPC.Jones.residual_lscAt} and \lean{HPC.Jones.exists_lscAt}.\\
Proposition~\ref{prop:lsc-homeo} & \lean{HPC.Jones.lscAt_homeomorph}.\\[4pt]
\multicolumn{2}{@{}l}{\emph{Chapter~\ref{ch:effros}}}\\
Proposition~\ref{prop:homeo-group} & \lean{HPC.Effros.homeoGroup}, \lean{HPC.Effros.homeoTopology}, \lean{HPC.Effros.homeo_isTopologicalGroup}, \lean{HPC.Effros.homeo_continuousSMul}, \lean{HPC.Effros.homeo_polishSpace}, \lean{HPC.Effros.supMove} (our $\mv$), \lean{HPC.Effros.supMove_inv}, \lean{HPC.Effros.supMove_mul_le}, \lean{HPC.Effros.hasBasis_supMove} (in \lean{HPC/Effros/HomeoGroup.lean}). Mathlib has no topology on the homeomorphism group at the pinned version; the formal proof uses the compact-open topology on $C(Y,Y)$, which for compact metric $Y$ is the topology of the uniform metric.\\
Definition~\ref{def:fat} & \lean{HPC.Effros.NowhereMeagre}, \lean{HPC.Effros.IsFat}.\\
Lemma~\ref{lem:fat-basic} & \lean{HPC.Effros.NowhereMeagre.diff_meagre} (van Mill's Lemma~2.1), \lean{HPC.Effros.NowhereMeagre.closure_subset_closure_diff}, \lean{HPC.Effros.nonempty_interior_closure_of_not_meagre}.\\
The nested construction (Lemma~\ref{lem:core-construction}) & \lean{HPC.Effros.CoreGood}, \lean{HPC.Effros.coreGood_step}, \lean{HPC.Effros.exists_eq_of_core}.\\
Reducing the parameter space (Lemma~\ref{lem:reduced}) & \lean{HPC.Effros.exists_reduced_param} (stated for $P=\N^\N$, which suffices since Mathlib defines analytic sets as continuous images of $\N^\N$).\\
van Mill, Proposition~2.2, relative form (Theorem~\ref{thm:prop22}) & \lean{HPC.Effros.inter_nonempty_of_nowhereMeagre}.\\
van Mill, Proposition~2.2 (Corollary~\ref{cor:fat-meet}) & \lean{HPC.Effros.inter_nonempty_of_isFat}. The hypothesis that the space be nonempty is needed: in the empty space the empty set is analytic and fat (\lean{HPC.Effros.isFat_empty_of_isEmpty}).\\
Lemma~\ref{lem:vm-basis} & \lean{HPC.Effros.exists_vanMill_basis}.\\
van Mill, Lemma~3.1 (Lemma~\ref{lem:vm31}) & \lean{HPC.Effros.vanMill_3_1}.\\
van Mill, Lemma~3.2 (Lemma~\ref{lem:vm32}) & \lean{HPC.Effros.vanMill_3_2}.\\
van Mill, Corollary~3.3 (Lemma~\ref{lem:vm33}) & \lean{HPC.Effros.vanMill_3_3}.\\
van Mill, Lemma~3.4 (Lemma~\ref{lem:vm34}) & \lean{HPC.Effros.vanMill_3_4} (as printed by van Mill, with the closure).\\
van Mill, Proposition~3.5 (Proposition~\ref{prop:vm35}) & \lean{HPC.Effros.vanMill_3_5} (as printed, with the closure).\\
The open mapping principle; Effros (Theorem~\ref{thm:effros-open}) & \lean{HPC.Effros.effros_open_mapping} and \lean{HPC.Effros.isOpenMap_smul_left}.\\
Effros, Ungar (Theorem~\ref{thm:ungar}) & \lean{HPC.Effros.effros_homogeneous}; for a compact subset $M$ of a metric space which is homogeneous in its own topology (the form used for plane continua), \lean{HPC.Effros.effros_homogeneous_set}.\\[4pt]

\end{longtable}
